\section{Flat pressure contact at a disappearing arithmetic peak}
\label{sec:v58-pressure-contact}
The drift bound of Theorem~\ref{thm:v57-damped-drift} makes the
moving Gaussian centers tight. It does not make them coincide on the
diffusive scale. There is a further piece of information at an exact
resonance: its covariance is the physical covariance. Keeping this
second-order contact in the pressure comparison gives a strictly
stronger conclusion. The argument is a scalar spectral principle;
it does not depend on a billiard norm once its physical inputs have
been verified.

\subsection{A visible-branch principle}
Let $\Lambda$ be a compact parameter space and let
$(M_\lambda,\nu_\lambda,T_\lambda)$ be probability-preserving systems.
Let $X_\lambda:M_\lambda\to\R^d$ be centered bounded observables,
with a common supremum bound. Write $S_mX_\lambda$ for their actual
orbit sums. The systems may be invertible or noninvertible. The
following hypotheses concern finitely many compact branch supports;
all neighborhoods and bounds may first be chosen locally and then
made uniform by a finite cover.

\paragraph{Pressure.}
On a real ball $|t|<t_0$ there is a real $C^3$ function
$\mathscr P_\lambda$ with uniform derivatives through order three,
such that
\begin{equation}\label{eq:v58-pressure-input}
 \int e^{-t\cdot S_mX_\lambda}\,d\nu_\lambda
       \le C e^{m\mathscr P_\lambda(t)},\qquad
 \mathscr P_\lambda(0)=0,\quad
 D\mathscr P_\lambda(0)=0,\quad
 D^2\mathscr P_\lambda(0)=\Sigma_\lambda.
\end{equation}
The real symmetric matrices $\Sigma_\lambda$ are continuous and
uniformly positive definite. This hypothesis is an exponential moment
of the original probability, not an order relation in a distribution
space.

\paragraph{Branches and physical visibility.}
For each branch there is a real frequency $a_j(\lambda)$ and a
centered holomorphic logarithm $\psi_{j,\lambda}$ near that frequency.
Its first three frequency derivatives are uniformly bounded and
continuous in the parameter. At the real maximum put
\begin{equation}\label{eq:v58-abstract-peak}
 \operatorname{Re}\psi_j(a_j)=-\kappa_j\le0,\quad
 D\psi_j(a_j)=iv_j,\quad
 -D^2\psi_j(a_j)=H_j,\qquad v_j\in\R^d.
\end{equation}
The peak data $a_j,\kappa_j,v_j,H_j$ are continuous on their compact
supports. The complex symmetric $H_j$ are bounded and have uniformly
positive real part. Near every point with $\kappa_j=0$ there are bounded
physical endpoint tests $a,d$ for which, on a fixed complex
neighborhood,
\begin{equation}\label{eq:v58-visibility-input}
 \begin{gathered}
 \int a(x)e^{iz\cdot S_mX_\lambda(x)}d(T_\lambda^m x)
                  \,d\nu_\lambda(x)
       =A_j^{a,d}(\lambda,z)e^{m\psi_j(z)}+E_{j,m}(\lambda,z),\\
 |A_j^{a,d}|\ge a_*>0,\qquad |E_{j,m}|\le C\rho^m.
 \end{gathered}
\end{equation}
where $\rho<1$. The tests may differ between neighborhoods. They
need not equal the endpoint tests defining a particular observation.
In particular its observation amplitude is allowed to vanish.

\paragraph{Second-order contact.}
On the compact zero-damping locus one has
\begin{equation}\label{eq:v58-contact-input}
                         H_j(\lambda)=\Sigma_\lambda.
\end{equation}
This is a separate spectral hypothesis. It is not inferred from
continuity of $v_j$ or from nonnegativity of the pressure alone.
Set
\[
 h_j=\|\Sigma_\lambda-\operatorname{Re}H_j\|,\qquad
 \omega(\delta)=\sup_{j,\lambda:\,\kappa_j\le\delta}
                         \|H_j-\Sigma_\lambda\|,
\]
where a supremum over an empty set is zero. Continuity, compactness
and \eqref{eq:v58-contact-input} imply $\omega(\delta)\to0$ as
$\delta\downarrow0$. No rate in $\delta$ is assumed.

\begin{theorem}[Drift at a flat pressure contact]
\label{thm:v58-flat-contact}
Under these hypotheses there are uniform constants $C,\kappa_0>0$
such that, for $0<\kappa_j\le\kappa_0$,
\begin{equation}\label{eq:v58-sharp-drift}
 |v_j|\le C\bigl(\sqrt{\kappa_jh_j}+\kappa_j^{2/3}\bigr),\qquad
 |v_j|^2\le C\kappa_j(h_j+\kappa_j^{1/3}).
\end{equation}
At $\kappa_j=0$ one has $v_j=0$. In particular
\begin{equation}\label{eq:v58-flat-ratio}
 \lim_{\delta\downarrow0}\sup_{0<\kappa_j\le\delta}
                  \frac{|v_j|^2}{\kappa_j}=0.
\end{equation}
For each fixed $L<\infty$,
$\sup_{m\kappa_j\le L}\sqrt m\,|v_j|\to0$ as $m\to\infty$.
The conclusions remain valid when a section or observation amplitude
vanishes at the resonance.
\end{theorem}
\begin{proof}
Take the physical witnesses in \eqref{eq:v58-visibility-input}.
For $z=a_j+it$ the modulus of their orbit pairing is at most
$\|a\|_\infty\|d\|_\infty\int e^{-t\cdot S_mX_\lambda}d\nu_\lambda$.
Jensen's inequality in \eqref{eq:v58-pressure-input} gives
$\mathscr P_\lambda(t)\ge0$ by taking $m$th roots. The same operation
in the witness expansion therefore proves
\[
 F_j(t):=\mathscr P_\lambda(t)
                  -\operatorname{Re}\psi_j(a_j+it)\ge0.
\]
The exponentially small remainder disappears in this step. The
amplitude lower bound is used before taking roots. Compactness of the
zero-damping locus gives finitely many witness neighborhoods with one
smaller tilt radius and common constants.

The frequency convention in \eqref{eq:v58-abstract-peak} gives
\begin{equation}\label{eq:v58-pressure-gap-taylor}
 F_j(t)=\kappa_j+v_j\cdot t
       +\tfrac12t^{\mathsf T}
                  (\Sigma_\lambda-\operatorname{Re}H_j)t+O(|t|^3).
\end{equation}
In particular the quadratic term is a \emph{difference} of
covariances. At zero damping, nonnegativity of $F_j$ and $F_j(0)=0$
give $v_j=0$.

For positive damping and $v_j\ne0$, insert
$t=-r v_j/|v_j|$ into \eqref{eq:v58-pressure-gap-taylor}. It follows
that
\begin{equation}\label{eq:v58-tilt-budget}
                    |v_j|\le\kappa_j/r+\tfrac12h_jr+C_3r^2.
\end{equation}
Choose
\[
 r=\min\{\kappa_j^{1/3},\sqrt{\kappa_j/h_j}\},
\]
with the second quantity interpreted as infinity when $h_j=0$.
For small $\kappa_j$ this lies in the fixed tilt ball. If
$h_j\le\kappa_j^{1/3}$, all three terms of
\eqref{eq:v58-tilt-budget} are bounded by $C\kappa_j^{2/3}$.
Otherwise $r=\sqrt{\kappa_j/h_j}$, and
$r^2=\kappa_j/h_j\le\sqrt{\kappa_jh_j}$; all three terms are
bounded by $C\sqrt{\kappa_jh_j}$. This proves the first bound, also
when $v_j=0$. Squaring proves the second.

All sufficiently small positive dampings lie in the finitely many
witness neighborhoods: otherwise a compact subsequence would limit
to an uncovered zero. Since $h_j\le\omega(\delta)$ when
$\kappa_j\le\delta$, \eqref{eq:v58-flat-ratio} follows. Finally
\[
 m|v_j|^2\le CL\{\omega(L/m)+(L/m)^{1/3}\}
              \quad\hbox{when }m\kappa_j\le L,
\]
which proves the diffusive assertion. The observation amplitude
has not entered any part of the proof.
\end{proof}

\subsection{A common Gaussian shape}
For accretive complex symmetric $H$ define
\[
 G_H^{(d)}(Z)=(2\pi)^{-d/2}(\det H)^{-1/2}
                         e^{-Z^{\mathsf T}H^{-1}Z/2},
                    \qquad Z\in\R^d.
\]
The determinant branch is the one continued from real positive
matrices on the convex accretive domain. Suppose the coefficients
$q_{j,m}(\lambda,y)$ obey
$|q_{j,m}|\le C e^{-m\kappa_j}$. Define
\[
 K_m(\lambda,y,Z)=\sum_jq_{j,m}(\lambda,y)
                         G_{H_j}^{(d)}(Z-\sqrt m\,v_j),\qquad
 S_m(\lambda,y)=\sum_jq_{j,m}(\lambda,y).
\]
This $S_m(\lambda,y)$ is a scalar coefficient, not an orbit sum.

\begin{corollary}[Uniform Gaussian shape through a resonance change]
\label{cor:v58-common-shape}
There are $a,C,\delta_0>0$ such that, for $m\ge2$ and
$0<\delta<\delta_0$,
\begin{align}
 &\sup_{\lambda,y,Z}e^{a|Z|^2}
 |K_m(\lambda,y,Z)-S_m(\lambda,y)g_{\Sigma_\lambda}(Z)|
 \notag\\[-1mm]
 &\hspace{17mm}\le
 C\{\omega(\delta)+\sqrt{\omega(\delta)}
                    +m^{-1/6}+e^{-m\delta/4}\}.
                                      \label{eq:v58-shape-modulus}
\end{align}
In particular its left side tends to zero. This is a statement about
the evaluated peak sum; it is not a density theorem for an arbitrary
physical source realizing the spectral hypotheses.
\end{corollary}
\begin{proof}
The weaker bound $|v_j|^2\le C_D\kappa_j$ holds on all compact
supports. Near zero it follows from Theorem~\ref{thm:v58-flat-contact}
and boundedness of $h_j$; away from zero use the positive minimum of
$\kappa_j$ and boundedness of the drift.

On the compact accretive matrix family and its line segments to
$\Sigma_\lambda$, Gaussian differentiation and completion of the
square give uniform bounds
\[
 |\nabla_ZG_H^{(d)}(Z)|+\|D_HG_H^{(d)}(Z)\|
                                  \le C e^{-b|Z|^2}
\]
for some $b>0$, after absorbing the polynomial factors. The same
bound holds for the Gaussian itself. Along the segment in the
matrix and in the shift $d_j=\sqrt m\,v_j$, the fundamental theorem
of calculus gives
\[
 |G_{H_j}^{(d)}(Z-d_j)-g_{\Sigma_\lambda}(Z)|
 \le C(\|H_j-\Sigma_\lambda\|+|d_j|)
             \sup_{0\le s\le1}e^{-b|Z-sd_j|^2}.
\]
Decrease $b$ so that $bC_D\le1/2$. Using
$|Z-sd_j|^2\ge|Z|^2/2-|d_j|^2$ shows that the product with
$e^{-m\kappa_j}$ is bounded by
\begin{equation}\label{eq:v58-shape-one-peak}
 C e^{-a|Z|^2}e^{-m\kappa_j/2}
                   (\|H_j-\Sigma_\lambda\|+\sqrt m\,|v_j|).
\end{equation}
For $\kappa_j\le\delta$, use
\eqref{eq:v58-sharp-drift}. The term
$\sqrt{m\kappa_jh_j}e^{-m\kappa_j/2}$ is at most
$C\sqrt{\omega(\delta)}$. The cubic-remainder term is bounded by
\[
 \sqrt m\,\kappa_j^{2/3}e^{-m\kappa_j/2}
       =m^{-1/6}(m\kappa_j)^{2/3}e^{-m\kappa_j/2}
       \le C m^{-1/6}.
\]
The Hessian difference costs $C\omega(\delta)$. For
$\kappa_j>\delta$, bounded Hessians and
$\sqrt m|v_j|\le C\sqrt{m\kappa_j}$ bound
\eqref{eq:v58-shape-one-peak} by
$Ce^{-a|Z|^2}e^{-m\delta/4}$. Sum the finite family. For instance,
$\delta=m^{-1/2}$ gives convergence. The unknown modulus $\omega$
prevents this choice from being advertised as a polynomial rate.
\end{proof}

\subsection{Verification on the occupation spectrum}
\begin{corollary}[Flat contact for the original Lorentz record]
\label{cor:v58-lorentz-contact}
For the actual occupation peaks of \eqref{eq:v41-peak-data},
\eqref{eq:v58-sharp-drift} holds with
$v_j=\mu_j-\mu_R$ and $\Sigma_R=\Omega_R$.
All constants are uniform for $R\in[0.45,0.47]$, including changes
of arithmetic type. No section residue is required to be nonzero.
\end{corollary}
\begin{proof}
Use $X_R=\mathsf f_R-\mu_R$, including the actual occupation
$\eta_R$, not the constant collision step. Lemma~\ref{lem:v57-centered-pressure}
supplies the pressure. Differentiation at zero of its full occupation
splitting gives $D^2\mathscr P_R(0)=\Omega_R$, by the same
Green--Kubo identification used in Section~\ref{sec:v37-occupation-local}.
Its third derivatives are uniformly bounded by strong frequency
analyticity on a fixed ball. Theorem~\ref{thm:v57-damped-drift}
constructs exactly the nonzero full-collision witness neighborhoods
in \eqref{eq:v58-visibility-input}, also where $B_j=0$.
Lemma~\ref{lem:v41-scalar-continuity} supplies scalar derivative
continuity, and Lemma~\ref{lem:v41-moving-maxima} supplies
$H_j=\Omega_R$ on $\kappa_j=0$ and the compact accretive bounds.
These verify each hypothesis separately. Apply
Theorem~\ref{thm:v58-flat-contact}. No parameter derivative, new
anisotropic space, or growing-band estimate is used.
\end{proof}

\subsection{An explicit singular hyperbolic realization}
The principle is not restricted to billiard collision operators.
Here is a deterministic example with a finite resonance transition,
for which all physical scalar inputs can be verified directly.
On the three strips $i/3\le x<(i+1)/3$, $i=0,1,2$, of the unit
square, define the baker transformation
\begin{equation}\label{eq:v58-baker-map}
 \mathcal B(x,y)=(3x-i,(y+i)/3),\qquad
 f_\epsilon(x,y)=\begin{cases}0&i=0,\\1&i=1,\\2+\epsilon&i=2.\end{cases}
\end{equation}
Boundary conventions affect a null set only. The map preserves
Lebesgue probability: each strip is mapped bijectively, with Jacobian
one, to a disjoint horizontal strip. Its derivative on each piece is
$\operatorname{diag}(3,1/3)$. It is an invertible singular hyperbolic
map, with a geometry different from a dispersing billiard.

\begin{proposition}[A damped arithmetic peak for the baker map]
\label{prop:v58-baker-realization}
For all sufficiently small real $\epsilon$, the centered observable
$f_\epsilon-\mu_\epsilon$, $\mu_\epsilon=1+\epsilon/3$, satisfies
the visible-branch and pressure hypotheses near the resonance $2\pi$.
Its physical covariance is
$\Sigma_\epsilon=2/3+2\epsilon/3+2\epsilon^2/9$.
The unique nearby real maximum $a_\epsilon$ has data
\begin{align}
 a_\epsilon&=2\pi-\pi\epsilon+\tfrac\pi3\epsilon^2+O(\epsilon^3),
                                      \label{eq:v58-baker-maximum}\\
 \kappa_\epsilon&=\tfrac{\pi^2}{9}\epsilon^2+O(\epsilon^3),\qquad
 v_\epsilon=\tfrac{\pi^2}{18}\epsilon^3+O(\epsilon^4),
                                      \label{eq:v58-baker-drift}\\
 H_\epsilon-\Sigma_\epsilon
             &=\tfrac{2\pi i}{9}\epsilon+O(\epsilon^2).
                                      \label{eq:v58-baker-hessian}
\end{align}
In particular $\epsilon_m\to0$ with $m\epsilon_m^2\to t<\infty$
gives $m\kappa_{\epsilon_m}\to\pi^2t/9$ but
$\sqrt m\,v_{\epsilon_m}\to0$. The peak can retain a nonzero
limiting damping factor without retaining a displaced Gaussian center.
\end{proposition}
\begin{proof}
Prescribing the first $m$ strip digits fixes an interval of $x$ of
length $3^{-m}$ and leaves $y$ arbitrary. Thus these digits are
independent and uniform under the invariant probability, as an exact
property of this deterministic map. Put
\[
 \varphi_\epsilon(z)=\tfrac13
                (1+e^{iz}+e^{i(2+\epsilon)z}).
\]
Its physical characteristic moment is exactly
$\varphi_\epsilon(z)^m$. The centered branch is
$e^{-iz\mu_\epsilon}\varphi_\epsilon(z)$; the witnesses are
$a=d=1$, their amplitude is one, and the remainder is zero.
The centered real exponential moment therefore has pressure
\[
 \mathscr P_\epsilon(t)=
           \log(e^{t\mu_\epsilon}\varphi_\epsilon(it)).
\]
Its first derivative vanishes and its Hessian is the stated variance.
The same identity proves the moment hypothesis at every finite count.
All complex derivatives are bounded on fixed neighborhoods of zero
and $2\pi$ for small $\epsilon$.

At $\epsilon=0$ the logarithmic real curvature at $2\pi$ is $-2/3$.
The analytic implicit function theorem gives a unique analytic nearby
maximum. The triangle equality $|\varphi_\epsilon(a)|=1$ requires
$e^{ia}=e^{ia(2+\epsilon)}=1$. On a small neighborhood of
$(\epsilon,a)=(0,2\pi)$ this holds only at that point. There the
centered branch is its origin branch multiplied by a constant phase;
hence $H_0=\Sigma_0$. This verifies the contact condition without
assuming it throughout the transition.

We give the expansion calculation. Write $a=2\pi+s$ and let
$F$ take the values $0,1,2+\epsilon$ with equal probabilities.
The three residual phases are
$\Theta=(0,s,2s+2\pi\epsilon+\epsilon s)$. For
$s=A\epsilon+B\epsilon^2+O(\epsilon^3)$, the first two coefficients
of $\operatorname{Re}(\varphi'/\varphi)$ are
\[
 -\tfrac23(A+\pi),\qquad
 -\tfrac23A-\tfrac23B-\tfrac49\pi.
\]
Their vanishing gives $A=-\pi$, $B=\pi/3$ and
\eqref{eq:v58-baker-maximum}. The real logarithm to second order is
$-\tfrac12\operatorname{Var}\Theta$, giving
$\kappa_\epsilon=\pi^2\epsilon^2/9+O(\epsilon^3)$.
For the imaginary logarithmic derivative, expansion of the finite
exponential sum gives
\[
 \operatorname{Im}(\varphi'/\varphi)-\mathbb EF
 =-\tfrac12\mathbb E[(\Theta-\mathbb E\Theta)^2
                                    (F-\mathbb EF)]+O(\epsilon^4).
\]
At the maximum the expectation on the right is
$-\pi^2\epsilon^3/9+O(\epsilon^4)$; its second-order term is zero.
Finally $-(\log\varphi)''$ equals
$\operatorname{Var}F+i\mathbb E[(F-\mathbb EF)^2
(\Theta-\mathbb E\Theta)]+O(\epsilon^2)$, and the expectation is
$2\pi\epsilon/9+O(\epsilon^2)$. These identities prove
\eqref{eq:v58-baker-drift}--\eqref{eq:v58-baker-hessian}. Analyticity
of the implicit maximum bounds the displayed remainders. The limiting
assertions follow by substitution.
\end{proof}

\paragraph{Scope of the second realization.}
This example verifies the same pressure/contact theorem in a piecewise
linear hyperbolic map by exact physical pairings, independently of the
billiard imports. Its additive observable has finite support, so no
continuous roof-density theorem is claimed for it. It is not substituted
for the original Lorentz record or its incidence and clearance sources.
